% 付録再編草案・第1部：更新を繰り返す

\section{更新を繰り返すと,何が見えるか}

本文で一次分数型漸化式を解いたとき,不動点を見つけ,変数を取り直すと,複雑な分数の更新が等比数列に変わった.
あの変数変換は,計算を進めるための工夫であると同時に,「見方を変えると,複雑な動きの中に単純な動きが隠れている」ことを教えてくれる.
この節では,その発見を出発点にして,漸化式の背後にある地形を少しずつ歩いてみよう.

まず,一項前だけで次が決まる漸化式
\[
  a_{n+1}=f(a_n)
\]
を考える.
ここで $a_n$ はその時点の状態,\,\,$f$ は状態を次の状態へ移す更新規則である.
本文の一次分数型では,変数を適切に取り直すことで,更新規則そのものを等比数列の形にできた.
では,変数を取り直しても全体を単純な式にできない場合はどうだろう.
数列がある値の近くを通っているときだけなら,動きを単純に見られないだろうか.

その値を $\alpha$ とし,そこからのずれを
\[
  e_n=a_n-\alpha
\]
と置く.
本文で見たように,\,\,$\alpha$ が不動点なら,\,\,$f(\alpha)=\alpha$ である.
数列が不動点の近くにいるときは,
\[
  e_{n+1}=f(\alpha+e_n)-f(\alpha)
\]
となる.
ここで微分係数を思い出そう.
グラフ上の点 $\alpha$ の近くでは,曲線をその点の接線で近似できる.
つまり,入力を少しだけ動かしたとき,出力がどれだけ動くかは傾き $f'(\alpha)$ で測れる.
この考えを式にすると,
\[
  f(\alpha+e_n)\approx f(\alpha)+f'(\alpha)e_n
\]
である.
したがって,
\[
  e_{n+1}\approx f'(\alpha)e_n
\]
となる.
不動点のすぐ近くでは,非線形な動きも「ずれを一定倍率で伸ばす・縮める」動きに見える.
傾きの絶対値が $1$ より小さければ,小さなずれは縮む.
本文の変数変換は規則全体を等比数列にしたが,ここでは接線が不動点の近くの動きを等比数列に見せている.

次に,過去の項を二つ使う漸化式へ進もう.
ここで行列に初めて出会う読者のために,行列が何をする道具なのかを確かめておく.
行列は数を四角く並べた表であると同時に,列ベクトルに掛けて新しい列ベクトルを作る規則を記録したものである.
たとえば
\[
  \begin{pmatrix}a&b\\c&d\end{pmatrix}
  \begin{pmatrix}x\\y\end{pmatrix}
  =
  \begin{pmatrix}ax+by\\cx+dy\end{pmatrix}
\]
では,一行目が新しい第一成分を,二行目が新しい第二成分を決めている.
各成分を別々に計算する規則を,一つの記号にまとめて持ち運べるわけだ.

二階線形漸化式
\[
  a_{n+2}=p a_{n+1}+q a_n
\]
では,いまの項 $a_{n+1}$ だけを見ても次の項は決まらない.
直前の項 $a_n$ も一緒に持っておく必要がある.
そこで $(a_{n+1},a_n)$ を一組の状態として,
\[
  \boldsymbol v_n=
  \begin{pmatrix}a_{n+1}\\a_n\end{pmatrix}
\]
と置く.
この状態に対し,
\[
  A=\begin{pmatrix}p&q\\1&0\end{pmatrix}
\]
を掛けると,
\[
  A\boldsymbol v_n
  =
  \begin{pmatrix}p a_{n+1}+q a_n\\a_{n+1}\end{pmatrix}
  =
  \begin{pmatrix}a_{n+2}\\a_{n+1}\end{pmatrix}
  =\boldsymbol v_{n+1}
\]
となる.
一行目はもとの漸化式そのもので,二行目は直前の第一成分を次の第二成分へ移している.
だから
\[
  \boldsymbol v_{n+1}=A\boldsymbol v_n
\]
は新しい仮定を加えた式ではなく,二つの項の更新を一つにまとめて書き直した式である.
同じ規則を二回使えば $\boldsymbol v_{n+2}=A^2\boldsymbol v_n$ となる.
ここで $A^2$ は $A$ を二度続けて掛けることを表す.
行列の累乗は,同じ更新を何歩も先まで繰り返す記号なのである.

では,行列を何度も掛けると何が起こるのか.
行列の更新は,ベクトルの足し算や定数倍と両立する.
この性質を線形性という.
そのため,更新しても向きが変わらず,倍率だけが変わる方向が見つかれば,何歩先の状態も簡単に追える.
その特別な方向を表す零でないベクトル $\boldsymbol v$ と倍率 $\lambda$ が
\[
  A\boldsymbol v=\lambda\boldsymbol v
\]
を満たすとき,\,$\boldsymbol v$ を固有ベクトル,\,$\lambda$ を固有値という.
この方向の状態は一歩ごとに
\[
  \boldsymbol v,\quad \lambda\boldsymbol v,\quad
  \lambda^2\boldsymbol v,\quad\ldots
\]
と変化する.
行列の反復が,その方向に限れば等比数列の反復になる.

この特別な方向がどの倍率 $\lambda$ に対して存在するかを調べよう.
$I$ を $I\boldsymbol v=\boldsymbol v$ となる単位行列とすると,
\[
  (A-\lambda I)\boldsymbol v=\boldsymbol 0
\]
である.
零でない $\boldsymbol v$ がこの式を満たすということは,\,$A-\lambda I$ がある方向をつぶして零にするということだ.
二次元では,これは行列の二つの列が同じ方向に並ぶことに当たる.
二つの列が
\[
  \begin{pmatrix}u\\w\end{pmatrix},
  \qquad
  \begin{pmatrix}v\\z\end{pmatrix}
\]
なら,同じ方向に並ぶ条件は $uz-vw=0$ である.
この $uz-vw$ を行列式と呼ぶ.
つまり,この行列式が $0$ になることが,零でない固有ベクトルを持つための条件である.
したがって,
\[
  \det\begin{pmatrix}p-\lambda&q\\1&-\lambda\end{pmatrix}=0
\]
となり,
\[
  \lambda^2-p\lambda-q=0
\]
を得る.
これは本文で使った特性方程式である.
特性方程式は,突然思いついた解法上の約束ではない.
二項をまとめた状態の中に,倍率だけで追える方向があるかを調べた結果だった.

二つの特性根 $\lambda_1,\lambda_2$ が異なるとき,それぞれの固有ベクトルは同じ直線上には並ばない.
もし同じ直線上なら,同じベクトルが $\lambda_1$ 倍と $\lambda_2$ 倍の両方になるため,二つの根が等しくなってしまうからだ.
平面では,向きの異なる二本のベクトルを使ってどんな状態も表せる.
つまり初めの状態は二つの固有ベクトルを定数倍して足した形で書け,\,\,$k$ 歩後にはそれぞれの成分が $\lambda_1^k$ 倍と $\lambda_2^k$ 倍される.
これが,一般項に二つの特性根の冪が現れる理由である.
根が複素数の場合も,いったん成分を複素数まで広げれば同じ計算ができる.
もとの漸化式と初期値が実数なら,二つの複素共役な成分が組になって,最終的な数列は実数になる.

フィボナッチ数列なら $p=q=1$ なので,行列は
\[
  A=\begin{pmatrix}1&1\\1&0\end{pmatrix}
\]
となり,固有値は $\varphi=(1+\sqrt5)/2$ と $\psi=(1-\sqrt5)/2$ である.
初めの状態 $(F_1,F_0)$ も,この二つの固有ベクトルの方向への成分を足した形で表される.
一歩進むごとに,状態の $\varphi$ 方向の成分は $\varphi$ 倍され,\,\,$ \psi$ 方向の成分は $\psi$ 倍される.
$|\psi|<1<\varphi$ だから,歩みを重ねるにつれて前者が全体の形を支配する.
一般項に現れる二つの冪は,この二つの方向の成長を一つの数列の上で見たものなのである.

重解のときは,二つの異なる倍率の方向に分けることができない.
では,一般項に現れる $n$ はどこから来るのだろう.
重なった特性根を $r\ne0$ とすると,特性方程式は $(\lambda-r)^2=0$ だから,漸化式は
\[
  a_{n+2}=2r a_{n+1}-r^2a_n
\]
となる.
ここで $a_n=r^n b_n$ と置いて代入すると,
\[
  b_{n+2}-2b_{n+1}+b_n=0
\]
を得る.
これは $b_n$ の隣り合う差が一定という式なので,\,\,$b_n=C+Dn$ である.
したがって
\[
  a_n=(C+Dn)r^n
\]
となり,重解のときに $nr^n$ が現れることが分かる.
たとえば $r=1$ なら,元の式は
\[
  a_{n+2}-2a_{n+1}+a_n=0
\]
である.
これを
\[
  (a_{n+2}-a_{n+1})-(a_{n+1}-a_n)=0
\]
と見れば,隣り合う項の差が一定で,\,$a_n=C+Dn$ となる.
なお $r=0$ の場合は $a_{n+2}=0$ となる退化した場合で,この漸化式を直接見ればよい.

固有値が複素数になると,数列はなぜ振動するのだろう.
複素数 $\rho(\cos\theta+i\sin\theta)$ は,複素平面で点を角 $\theta$ だけ回し,長さを $\rho$ 倍する操作を表す.
その累乗は角を $n\theta$ にし,長さを $\rho^n$ にする.
したがって,複素共役な特性根からは,振幅が増減しながら振動する数列が生まれる.

この振動を,多項式として保存する例がチェビシェフ多項式である.
\[
  T_0(x)=1,\quad T_1(x)=x,\qquad
  T_{n+1}(x)=2xT_n(x)-T_{n-1}(x)
\]
と定める.
$x$ を固定すると,これは二階線形漸化式で,特性方程式は
\[
  r^2-2xr+1=0
\]
である.
$x=\cos\theta$ と置くと,特性根は $\cos\theta\pm i\sin\theta$ となる.
一方,余弦の加法定理も
\[
  \cos((n+1)\theta)
  =2\cos\theta\cos(n\theta)-\cos((n-1)\theta)
\]
という同じ更新規則を持つ.
初めの値も一致するため,
\[
  T_n(\cos\theta)=\cos(n\theta)
\]
である.
たとえば $T_2(x)=2x^2-1$ は $\cos(2\theta)$ を $\cos\theta$ の多項式に翻訳し,\,\,$T_3(x)=4x^3-3x$ は三倍角を同じように翻訳する.
漸化式は数列の項だけでなく,多項式そのものも次々に作っていた.
$-1\le x\le1$ では $x=\cos\theta$ と表せるため,\,\,$T_n(x)$ は振動する.
一方,\,\,$x=\cosh t>1$ と置けば
\[
  T_n(\cosh t)=\cosh(nt)
\]
となり,同じ多項式列が今度は急速な増大を表す.
境界の $x=1$ では $T_n(1)=1$ であり,振動と増大の切り替わりが一つの式の中に収まっている.

ここからもう一歩だけ景色を広げよう.
余弦関数どうしは,区間 $[0,\pi]$ で積分すると,次数が異なるものどうしの面積が打ち消し合う.
この性質を直交性という.
$x=\cos\theta$ と変数を取り直すと,チェビシェフ多項式について
\[
  \int_{-1}^{1}
  \frac{T_m(x)T_n(x)}{\sqrt{1-x^2}}\,dx=0
  \qquad(m\ne n)
\]
となる.
直交する多項式は,ベクトルの直交する方向のように,関数を成分へ分ける基準として使える.
そのためチェビシェフ多項式は,三角関数の公式にとどまらず,関数近似や数値計算にも現れる.
特性方程式から始めた道が,回転・振動・多項式・近似へ伸びている.
